\section{Introduction}

Trust in media outlets in Mexico reached a historic low in 2024, continuing a downward trend that began in 2017 \citep{Reuters}. This phenomenon is not unique to Mexico; globally, trust in media has been declining, with trust in the U.S., for example, falling to 31\% \citep{Gallup-US}. In democratic contexts, this decline is especially concerning, as a well-functioning information environment plays a critical role in holding governments accountable, informing the public, and fostering transparency and responsiveness \citep{Stromberg-MediaCoverage, Account-Boix-JoLEO}. In Latin America, the implications are even more troubling due to lower state capacity, weaker institutions \citep{Audits-Ferraz-QJE}, and violent contexts \citep{Marshall}, all of which undermine the availability of obtaining accurate information which is essential for democratic accountability.

The economic and political science literature, as well as reports from journalistic organizations such as the \href{https://reutersinstitute.politics.ox.ac.uk/}{Reuters Institute}, often highlight external factors to explain this decline. These include political leaders delegitimizing media outlets \citep{PressCriticism-Carlson-Journalism}, the rise of populism \citep{Populism-MediaTrust, PolEconPop-Guriev-JEL}, the growth of social media, and increased political polarization \citep{Account-Larreguy-WP}. However, relatively few studies have examined the role of the media itself, including its incentives \citep{ArgCorr-DiTella-AEJAE}, the perceived bias of outlets \citep{Mullainathan-MediaBias, Stromberg-MediaBias, Gentz-MediaBias}, and how these factors influence their behavior and decisions.

In Mexico, the decisions media outlets make—what stories to cover, how to cover them, and when—have been deeply tied to the government’s use of state resources to influence narratives, suppress dissent, and maintain political dominance. From outright authoritarian control during the 20th century \citep{FourthState-Lawson-Book} to subtler forms of influence through advertising spending \citep{Times}, the relationship between the media and the state remains a key factor in shaping public discourse. This study seeks to examine how a significant change in an external factor, such as the government’s financial support of media outlets through advertising spending, affects critical coverage of the government. Specifically, I focus on the abrupt reduction in advertising spending between two administrations and its effect on media behavior.

Under the \textit{republican austerity} policy implemented during Andrés Manuel López Obrador’s (AMLO) presidency, the government reduced its advertising spending by 78.7\%, a sharp decrease that likely influenced media coverage of certain issues \citep{Article19}. I hypothesize that this substantial reduction led to an increase in critical coverage of the government, particularly on the salient issue of violence. In 2018, a year before these cuts, Mexico recorded its highest-ever intentional homicide rate, with 30 homicides per 100,000 inhabitants annually under the Peña Nieto administration \citep{Intentional-Homicides}.

Several mechanisms could be at play. Based on the work of \cite{Censor-Pratt-AER}, the reduction in advertising spending may have weakened the government’s ability to control the information environment, allowing media outlets greater freedom to report on the actual levels of violence during the first two years of AMLO in office compared to the last two years of Peña Nieto. Another possible explanation is that media outlets, responding to the funding cuts, adopted a more critical stance either as retaliation or as a strategy to pressure the government into reinstating higher budgets. Media owners, feeling the financial strain of reduced funding, may have encouraged negative coverage as a way to influence public opinion and the administration’s policies. Alternatively, this reaction could reflect a broader effort of the media to restore trust in their reporting.

I will also focus on the topic of violence, as it remained a consistently salient issue across successive administrations. This contrasts with topics like corruption, which are more time-sensitive and often depend on how long officials have been in office. It would be unusual for newly appointed officials to commit acts of corruption immediately after taking office, and it would also be challenging for journalists to produce credible stories about corruption within such a short timeframe. Other issues I could focus on could potentially be the economy and foreign affairs, like the relationship with the Trump administration, which also remained salient after the change.

To study the causal effect of the reduction in advertising spending, I will use a difference in difference (DiD) design with continuous treatment, using media outlets as units, and seeing their coverage across periods and across groups. The continuous treatment variable will be the percentage increase or decrease in government spending after the administration change. Given that violence levels could be different across periods, both in magnitude and distribution, I will also conduct an analysis at the state level, assigning each local media outlet to their respective state and including state-time fixed effects, to leverage the intra-state variation in violence coverage.

For the analysis, I have collected data on government advertising spending (COMSOC), which includes the date (at the daily level) each contract was enforced, the name of the media outlet receiving funds, the name of the government agency contracting the media services or purchasing airtime, and the amount of money allocated. To compile a comprehensive dataset of newspaper articles that can be matched to this spending data, I will scrape a near-complete set of articles published online by newspapers from 2012 to 2024. I will use web crawlers for data collection, then employ ML and large language models (LLMs) to label each article as political, violence-related, positive, or negative.

To the best of my knowledge, this would be the first study to causally link a government censorship attempt on media behavior in a democratic context. Some anecdotal evidence has been documented on the use of resources to silence the press \citep{Peru-Macmillan-JEP, ArgCorr-DiTella-AEJAE}, as well as other legal methods that hinder the access to sources or the ability to write articles \citep{Kisha-UoNott, Kisha-Freedom, MexicanPressStanigAJPS}. In authoritarian contexts, most of the literature focuses on China \citep{CensorshipChina-King-Science, CensorshipCollective-King-APSR, Army-Chen-Book, Plus-Anne-PostComm, PersonnelMedia-Lee-IntJPPoli} while also setting the theoretical basis for other contexts \citep{PoorDict-Guriev-APSR, ElitesMedia-Tung-PolComm, Autocracy-Guriev-JPE, Censor-Pratt-AER}.

The rest of the research proposal is organized as follows. Section \ref{sec:lit_review} provides a detailed description of the literature related to this topic and the historical and political Mexican context. Section \ref{sec:data} outlines a description of all the data used in the analysis, as well as an explanation on how I will build each variable. Section \ref{sec:emp_strat} presents the empirical strategies, as well as discussion around them. Section \ref{sec:results} presents preliminary results.

